\chapter{Local gradient bounds from height oscillation}\label{sec:graphical-height-gradient}
An ambient-ball cutoff controls angle only while the surface remains inside that ambient ball. A height-dependent cutoff instead gives a gradient estimate on a fixed base ball. We develop the maximum-principle argument of \cite{EH91}, Theorem 2.3, printed pp.~551--553, with explicit, nonoptimal constants. The flow is assumed smooth on the relevant closed cylinder; this section proves an estimate for such a flow, not the boundary-value existence or approximation theorem.

\begin{lemma}[The height-dependent cutoff at a large-angle maximum]\label{lem:graphical-height-cutoff-gradient}
Work on the unit base ball, suppose the graph height $h<0$ throughout a finite time slab, and let $\beta\geq1$. On the positive set of
\[
 \varphi=1-|x|^2+\frac{h}{2\beta}
\]
one has $\mathcal P\varphi\leq2n$. At a point with $v\geq4+16\beta$, one also has
\[
 |\nabla\varphi|^2\geq\frac1{16\beta^2}.
\]
Here $\mathcal P=\partial_t^\perp-\Delta$ and gradients are intrinsic to the moving graph.
\end{lemma}
\begin{proof}
The ambient height satisfies $\mathcal Ph=0$. Lemma \ref{lem:graphical-compact-cutoff-geometry} gives $\mathcal P(1-|x|^2)=2\operatorname{tr}(g^{-1})\leq2n$. Write $e$ for the upward vertical unit vector and $X=(x,h)$. The intrinsic height gradient is $e^T$, so $|\nabla h|^2=1-v^{-2}$. Since the ambient gradient of $|x|^2$ is the horizontal vector $2(x,0)$,
\[
 \langle\nabla h,\nabla|x|^2\rangle
       =-2v^{-1}\langle\nu,(x,0)\rangle,
 \qquad
 |\langle\nabla h,\nabla|x|^2\rangle|\leq2/v
\]
in the unit base ball. Consequently
\[
 |\nabla\varphi|^2
 \geq\frac{1-v^{-2}}{4\beta^2}-\frac2{\beta v}.
\]
For $v\geq4+16\beta$ the first term is at least $3/(16\beta^2)$ and the subtracted term at most $1/(8\beta^2)$. This proves the estimate. Dropping the nonnegative term $|\nabla|x|^2|^2$ is legitimate; no estimate of the Euclidean height gradient by one is being made.
\end{proof}

\begin{theorem}[An explicit local height-to-gradient estimate]\label{thm:graphical-height-gradient-estimate}
Let a graphical flow be smooth on a neighborhood of $\overline B_R(y)\times[0,T]$, $R>0$. Define
\[
 V_0=\sup_{B_R(y)}\sqrt{1+|Du(x,0)|^2},\quad
 M=\sup_{B_R(y)\times[0,T]}u,\quad
 K=M-\inf_{0\leq t\leq T}u(y,t).
\]
Then $K\geq0$ and, for every $0\leq t\leq T$,
\[
 v(y,t)\leq42V_0\exp\left((32n+1)\max\{1,K^2/R^2\}\right).
\]
This is an explicit version of \cite{EH91}, Theorem 2.3. The constants here are sufficient constants, not sharp ones. Only the initial slope in the base ball and a height oscillation over the cylinder enter the estimate.
\end{theorem}
\begin{proof}
Translate and parabolically rescale to $y=0$, $R=1$, scaling height by $R^{-1}$ as well. First subtract $M+\epsilon$ from the rescaled height, so that the resulting $h$ is strictly negative throughout the cylinder. Set
\[
 \beta=\max\{1,K+\epsilon\},\qquad
 \lambda=64n\beta^2,\qquad
 \eta=(e^{\lambda\varphi}-1)_+,
 \quad\varphi=1-|x|^2+h/(2\beta).
\]
The positive support of $\eta$ stays strictly inside the base ball, uniformly over the finite slab. At the center $\varphi(0,t)\geq1/2$, whereas everywhere $\varphi<1$.

Consider a positive maximum of $v\eta$ over space and all times up to a chosen time. If it exceeds the initial supremum, it occurs at positive time and inside the positive support. There $\nabla(v\eta)=0$. The angle evolution from Proposition \ref{prop:graphical-angle-evolution} and the product rule give
\[
 0\leq\mathcal P(v\eta)
      =v\mathcal P\eta-|A|^2v\eta
      \leq v\mathcal P\eta.
\]
The two gradient terms cancel at the maximum. The tangential coordinate drift also vanishes there. Since $\eta=e^{\lambda\varphi}-1$ in this positive region, it follows that
\[
 \mathcal P\varphi\geq\lambda|\nabla\varphi|^2.
\]
If $v\geq4+16\beta$, Lemma \ref{lem:graphical-height-cutoff-gradient} would give $2n\geq4n$, a contradiction. Thus at such a maximum $v<4+16\beta\leq20\beta$. Since $\eta\leq e^\lambda$, the complete maximum estimate is
\[
 v\eta\leq e^\lambda\max\{V_0,20\beta\}
          \leq e^\lambda(V_0+20\beta).
\]
At the center $\eta\geq e^{\lambda/2}-1\geq\tfrac12 e^{\lambda/2}$, because $\lambda\geq64$. It follows that
\[
 v(0,t)\leq2(V_0+20\beta)e^{32n\beta^2}
       \leq42V_0e^{(32n+1)\beta^2},
\]
using $V_0\geq1$, $\beta\geq1$ and $\beta\leq e^{\beta^2}$. Let $\epsilon\downarrow0$ and restore the scale. In the rescaled problem $K$ is the original oscillation divided by $R$, giving the stated formula. Only the positive region of $\eta$ was differentiated, and the lateral zero set cannot carry a positive maximum.
\end{proof}

\begin{proposition}[Local a priori slope and curvature from a larger initial ball]\label{prop:graphical-initial-local-geometry-control}
Let an entire graphical flow be smooth on $[0,T]$. Fix $r>0$ and $R$ with $R^2>16r^2+2nT$. Suppose $|u(x,0)|\leq M$ on $B_R(y)$ and $v(x,0)\leq V_0$ on $B_{3r}(y)$. Define
\[
 V=42V_0\exp\left((32n+1)
          \max\{1,4(M+R)^2/r^2\}\right).
\]
Then $v\leq V$ on $B_{2r}(y)\times[0,T]$, and
\[
 \sup_{B_r(y)}|A(\cdot,t)|^2
       \leq C(n)V^4(t^{-1}+r^{-2})\qquad(0<t\leq T).
\]
The input is entirely local initial height and slope data. The height bound may be improved by subtracting a constant before applying the assertion.
\end{proposition}
\begin{proof}
Proposition \ref{prop:entire-graph-local-height-barrier} gives $|u|\leq M+R$ on $B_{4r}(y)\times[0,T]$. For any center $z\in B_{2r}(y)$, apply Theorem \ref{thm:graphical-height-gradient-estimate} to $B_r(z)$, which is contained in $B_{3r}(y)$. Its initial slope supremum is at most $V_0$, and its height oscillation $K$ is at most $2(M+R)$. Thus its center slope is bounded by $V$, uniformly in $z$ and time. Next apply Theorem \ref{thm:graphical-base-interior-curvature} with base radius $2r$ and interior ratio $1/2$. Absorb its fixed ratio factor in $C(n)$. Smoothness on finite compact cylinders justifies each maximum argument. Passing this estimate to merely locally Lipschitz initial data still requires the approximation/existence input in Theorem \ref{thm:ecker-huisken-arbitrary-growth-existence}; it is not inferred by differentiating that initial graph.
\end{proof}
